\section{The Regime Map: Which Curator to Reach For}
\label{sec:regime}

The experiments refuse a one-line verdict, and that refusal is the second
contribution. Outcome-based selection is not universally best; it is
best in a specific, characterizable regime, and worse outside it. Stating
that regime precisely is more useful---and more defensible---than a claim
that the ledger dominates, which our own tables would falsify. The axes
are ordered by how much they decide: the first two settle the question on
their own, and the rest are refinements.

\paragraph{Axis 1: Is there a stable, observable outcome?}
This is the gating precondition, not a tuning choice. The mechanism
requires a stable outcome definition and a known reporting path.
Passing tests and resolved issues are outcomes, not conserved quantities. Where
the only signal is a model scoring an answer---chat preference,
open-ended QA quality, ``helpfulness''---nothing pushes back, upkeep
would starve the long tail of correct-but-rarely-used knowledge, and the
mechanism is the wrong tool by construction. Systems like Mem0
\citep{mem0}, Zep \citep{zep}, Letta \citep{letta} and A-MEM
\citep{xu2025amem} serve that market and \textsc{darwin-memo}
deliberately does not. The benchmarks it is measured on---LoCoMo
\citep{maharana2024locomo}, LongMemEval \citep{wu2024longmemeval} and
their successors---score retrieval quality against reference answers, and
a store whose entries must \emph{earn} energy has nothing to earn from
there. We report no numbers on them for that reason, not because the
comparison is unfavourable.

\paragraph{Axis 2: Do measurements lie \emph{on purpose}?}
This is the axis the paper is about, and it decides on its own. What
separates the mechanisms under a curation-targeted adversary
(\S\ref{sec:adversary}) is not how aggressively they curate but what a
lie costs: a counter spends a life it cannot recover, a bounded credit
buffer spends energy an entry can earn back. The buffer stops being
ornamental the moment the noise has an author.

Three cautions keep this from being a victory lap. The bandit and
\texttt{keep\_everything} hold capability under every budget we test only
by never removing anything, so they are abstaining rather than defending.
The ledger has its own boundary, between two and four lies per cycle,
which we publish because a defence with an unstated limit is one nobody
can deploy safely. And our counter numbers are four settings, not a
family: sweeping the dial (\S\ref{sec:countersweep}) refuted our own
prediction, and a forgiveness counter at $k{=}5$ or $k{=}8$ holds both
axes through two lies per cycle. So the instruction here is narrower than
``use the ledger''. If your store may grow without cost and you are
willing to tune $k$, a consecutive-strike counter is a legitimate choice
under a small attack budget, and it is one line of code. Reach for the
ledger when the store's size is itself a cost, when you cannot bound the
adversary's budget in advance, or when the counter you already have is
the lifetime-strike shape, which has no working setting at any $k$ we
tested.

\pointer{Four axes refine rather than decide---what the adversary is
buying, accidental lying, whether refusals are measured, and whether the
horizon outruns the upkeep time constant, the condition our real-task leg
failed. \S\ref{sec:axes} states each in full.}
\inplace{regime-axes}

\paragraph{Summary.}
Outcome-based selection is the right curator when (1) a stable
outcome definition exists, (2) measurements are noisy in a one-sided way or
corrupted by an adversary, (3) inaction is priced or paired with
recovery, (4) the deployment runs for several multiples of the upkeep
time constant, and (5) leanness, complete removal of harmful entries,
label-free operation, or near-zero settlement cost are required. Outside
that region a one-line heuristic, a confidence bandit or a quarantine
policy can match or beat it, and we have shown exactly where---including
on the one real-task benchmark we ran, where condition (4) failed. The reporting path remains a trust boundary in every regime; no
retention rule makes a corrupted report truthful. Deployment cost and
success preservation require separate controlled evidence.
